\section*{第 \ref{ch:g12:buffers-predominance} 章——缓冲溶液与优势区域图}

\begin{solution}{exo:g12:buffers-predominance:1}
pH 2：\ce{CH3COOH}。pH 4.76：两者等量。pH 7：\ce{CH3COO-}。
\end{solution}

\begin{solution}{exo:g12:buffers-predominance:2}
pH 4：$1/(1 + 10^{0.76}) = 0.15$。pH 6：$1/(1 + 10^{-1.24}) = 0.95$。
\end{solution}

\begin{solution}{exo:g12:buffers-predominance:3}
黄色；绿色（在其变色范围内）；蓝色。
\end{solution}

\begin{solution}{exo:g12:buffers-predominance:4}
pH 7.4 时（低于 9.26）为 \ce{NH4+}。pH 11 时为 \ce{NH3}。
\end{solution}

\begin{solution}{exo:g12:buffers-predominance:5}
酸碱指示剂本身就是一对酸碱共轭对：
用量较多时会改变它原本要指示的 pH。
\end{solution}

\begin{solution}{exo:g12:buffers-predominance:6}
$10^{\mathrm{pH} - 4.76}$：0.1；1；10。
\end{solution}

\begin{solution}{exo:g12:buffers-predominance:7}
酚酞（8.0--10.0）或百里酚蓝（8.0--9.1）：它们的变色范围
均包含 8.7。甲基红变色于 4.2 至 6.3 之间，远离 8.7：它
早在远未到达 8.7 时就已变色。
\end{solution}

\begin{solution}{exo:g12:buffers-predominance:8}
$\mathrm{pH} = 4.76 + \log(0.10/0.20) = 4.76 - 0.30 = 4.46$.
\end{solution}

\begin{solution}{exo:g12:buffers-predominance:9}
pH 1：阳离子 \ce{H3N+-CH2-COOH}，电荷为 $+1$。pH 6：两性离子，
总电荷为 0。pH 12：阴离子 \ce{H2N-CH2-COO-}，电荷为 $-1$。
\end{solution}

\begin{solution}{exo:g12:buffers-predominance:10}
水：由 7.0 降至 $-\log(\num{1.0e-4}/0.101) = 3.0$，下降了 4 个 pH 单位。
缓冲溶液：由 4.76 降至 $4.76 + \log(9.9/10.1) = 4.75$，下降 0.01 个 pH 单位。
\end{solution}

\begin{solution}{exo:g12:buffers-predominance:11}
\ce{NH4+}/\ce{NH3}（$pK_a = 9.26$）。
比值 $[\ce{NH3}]/[\ce{NH4+}] = 10^{9.5 - 9.26} = 1.7$。
\end{solution}

\begin{solution}{exo:g12:buffers-predominance:12}
当比值达到 $1/10$ 时，pH 下降一个单位：
$(10 - n)/(10 + n) = 0.1$，故 $n = \qty{8.2}{mmol}$。
两种形式的量越多的缓冲溶液，在比值变化相同时
可吸收更多的酸或碱，即缓冲容量更大。
\end{solution}

\begin{solution}{exo:g12:buffers-predominance:13}
比值 $10^{5.0 - 4.76} = 1.74$；酸的量为 \qty{10}{mmol}，故
乙酸根 \qty{17.4}{mmol}，即 \qty{174}{mL} 的该溶液。
\end{solution}

\begin{solution}{exo:g12:buffers-predominance:14}
它含有两对酸碱共轭对，具有两个 $pK_a$，因而会发生两次
颜色变化：pH 1 时为红色，pH 5 时为黄色，pH 10 时为蓝色。
\end{solution}

\begin{solution}{exo:g12:buffers-predominance:15}
稀释使两种浓度均除以 10，比值
不变，故 pH（4.76）亦不变。强酸稀释十倍后，pH 升高
一个 pH 单位。
\end{solution}

\begin{solution}{pb:g12:buffers-predominance:1}
\textbf{1.} \ce{CO2 + 2H2O <=> HCO3- + H3O+}.

\textbf{2.} 溶解的二氧化碳（与水结合）是酸，
碳酸氢根是碱。

\textbf{3.} 表中数值是 \qty{25}{\celsius} 时纯水中的值；
血液温度为 \qty{37}{\celsius}，且含有大量其他离子，这些离子会改变
平衡常数（生理学家计入全部溶解的二氧化碳）。

\textbf{4.} 浓度为 \qty{24}{mmol/L} 时，$0.024 \times 5.0 = \qty{0.12}{mol}$。

\textbf{5.} pH 轴在 6.1 处分界：低于 6.1 为 \ce{CO2}，高于 6.1 为 \ce{HCO3-}。

\textbf{6.} pH 为 7.4 时，高于 6.1，为碳酸氢根。

\textbf{7.} 溶液略呈碱性：其 pH 大于 7。

\textbf{8.} $10^{7.4 - 6.1} = 10^{1.3} = 20$.

\textbf{9.} $24 / 20 = \qty{1.2}{mmol/L}$.

\textbf{10.} $10^{1.28} = 19$ 和 $10^{1.32} = 21$。

\textbf{11.} $20/21 = \qty{95}{\percent}$.

\textbf{12.} $10^{7.6 - 6.1} = 10^{1.5} = 32$.

\textbf{13.} 碳酸氢根离子：
\ce{HCO3- + H3O+ -> CO2 + 2H2O}.

\textbf{14.} $10^{7.1 - 6.1} = 10$.

\textbf{15.} 由 $24/20 = 1.2$ 增至 $24/10 = \qty{2.4}{mmol/L}$，即增加一倍。
肺部加快并加深呼吸，以呼出多余的
二氧化碳。

\textbf{16.} 弱酸及其共轭碱各自都能吸收所加入的酸或碱，
而 pH 仅取决于二者的比值；强酸则只会使
pH 朝一个方向变化。

\textbf{17.} 不会：比值与 pH 均保持不变。

\textbf{18.} 约为 20。
\end{solution}
